
\question
Galvanotaxis occurs when 

\begin{choices}
\choice plants respond to gravity.  
\choice protists respond to touch. 
\correctchoice protists respond to an electrical field. 
\choice animals respond to temperature. 
\choice plants respond to hormones. 
\choice animals respond to metal ions. 
\end{choices}

\question
Most of the water taken up by a plant

\begin{choices}
\choice is converted to sugars during photosynthesis. 
\choice transports sugars in the phloem. 
\correctchoice is lost during transpiration. 
\choice keep guard cells turgid. 
\choice enters the central vacuoles of cells.  
\end{choices}

\question
In terms of structure, cilia  are most like 
\begin{choices}
\choice myosin filaments. 
\choice bacterial flagella. 
\correctchoice eukaryotic flagella. 
\choice actin filaments. 
\choice microtubules.  
\end{choices}

\question
When ATP attaches to the myosin head

\begin{choices}
\choice the actin head attaches to the myosin filament. 
\choice the myosin head attaches to the actin filament.  
\correctchoice the myosin head releases from the actin filament. 
\choice the actin head attaches to the myosin filament. 
\choice the myosin head initiates the power stroke. 
\end{choices}

\newpage

\question
The fluid that moves around in the circulatory system of a typical arthropod is the

\begin{choices}
\choice blood plasma. 
\correctchoice hemolymph. 
\choice cytoplasm. 
\choice gastrovascular fluid. 
\choice intracellular fluid. 
\end{choices}

\question
The highest quality form of energy (best able to do work) for organisms is

\begin{choices}
\choice kinetic.
\correctchoice radiant. 
\choice potential.
\choice thermal.
\choice chemical. 
\end{choices}

\question
During protist conjugation, which of the following does \emph{not} occur during some part of the process?

\begin{choices}
\choice The micronucleus undergoes mitosis. 
\choice Some micronuclei distintegrate. 
\choice The macronucleus distintegrates. 
\choice The micronucleus undergoes meiosis. 
\choice Cells exchange micronuclei.  
\correctchoice The macronucleus fuses with the micronucleus. 
\end{choices}

\question
The ion used most often as a second messenger is

\begin{choices}
\correctchoice \ce{Ca+} 
\choice \ce{K+} 
\choice \ce{Cl-} 
\choice \ce{Na+} 
\choice \ce{H+} 
\end{choices}

\question
The main stimulus used by birds to find their way north for the summer time is 
\begin{choices}
\correctchoice electomagnetic.
\choice mechanical. 
\choice thermal. 
\choice chemical. 
\choice thigmic. 
\end{choices}

\question
Transpiration in plants requires all of the following \emph{except}

\begin{choices}
\choice adhesion of water molecules to cellulose. 
\choice evaporation of water molecules. 
\choice cohesion between water molecules.  
\choice water vapor on the spongy mesophyll. 
\correctchoice active transport through xylem cells. 
\end{choices}

\question
Differences in energy quality and availability affects the amount of surface area and volume in organisms. Compared to animals, plants have 

\begin{choices}
\correctchoice high surface area and low volumes. 
\choice low surface area and high volumes. 
\choice low surface areas and low volumes.  
\choice high surface area and high volumes. 
\end{choices}

\question
The function of valves in veins in a closed circulatory system is to 
\begin{choices}
\choice provide support for the extra muscles around the vein. 
\choice decrease the amount of pressure in the veins.  
\correctchoice ensure one way flow of circulatory fluids. 
\choice enable the vein to withstand high pressure. 
\choice enhance diffusion of nutrients to the tissues. 
\end{choices}

\question
Pores on the leaf surface that function in gas and water exchange are called

\begin{choices}
\choice phloem. 
\choice xylem. 
\choice waxy cuticles. 
\choice spongy mesophyll. 
\correctchoice stomata. 
\end{choices}


%\newpage

\question
Hydrostatic skeletons are 

\begin{choices}
\choice external skeletons that provide support and protection against dehydration. 
\correctchoice fluid-filled compartments under pressure. 
\choice internal skeletons that are buried in soft tissue.  
\choice skeletons made of microtubules in the cytoplasm of protists. 
\choice skeletons found in deep sea organisms that depend on high pressure of the surrounding water to provide support. 
\end{choices}

\question
What drives the flow of water through the xylem?

\begin{choices}
\choice Transpirational pull through sieve-tube elements. 
\choice Abscisic acid opening guard cells around the stomata. 
\choice A higher concentration of sugars in the leaves than in the roots. 
\correctchoice The evaporation of water from the leaves. 
\choice The number of companion cells in the phloem.  
\end{choices}

\question
The type of nerve cell responsible for integrating different stimuli is the 

\begin{choices}
\choice motor neuron 
\choice axon 
\choice dendrite 
\correctchoice interneuron 
\choice sensory neuron 
\end{choices}

\question
The protein that is activated in response to auxin that allows phototropism to occur is

\begin{choices}
\correctchoice expansin. 
\choice proton channel. 
\choice wallulase. 
\choice transmembrane. 
\choice abscisic acid. 
\choice basal apparatus. 
\end{choices}

\question
Organisms with a circulating fluid that is separate from the fluid that directly surrounding the body's cells are likely to have

\begin{choices}
\choice a vascular system. 
\choice a hemolymph system. 
\choice a gastrovascular cavity. 
\correctchoice a closed circulatory system. 
\choice an open circulatory system. 
\end{choices}

\question
From an energy perspective, animals are mobile because they consume

\begin{choices}
\choice high quality energy and carbon intake is irregular. 
\choice low quality energy and carbon intake is constant. 
\choice high quality energy and carbon intake is constant. 
\choice low quality energy and low energy efficiency.
\correctchoice low quality energy and carbon intake is irregular.
\choice high quality energy and high energy efficiency. 
\end{choices}

\question
The vascular tissue system for delivery of water, minerals, and nutrients throughout the plant includes all of the following except for

\begin{choices}
\correctchoice cohesion. 
\choice sieve plates. 
\choice phloem. 
\choice companion cells. 
\choice xylem. 
\end{choices}

\question
Ignoring all other factors, what kind of day would result in the fastest delivery of water and minerals to the leaves of a tree?

\begin{choices}
\correctchoice A warm, dry day. 
\choice A cool, humid day. 
\choice A cool, dry day. 
\choice A very hot, dry, windy day. 
\choice A warm, humid day.  
\end{choices}

\newpage

\question
Which of the following are primarily responsible for movement in the eukaryotic flagellum?

\begin{choices}
\choice Actin and myosin filaments.  
\choice Hydrogen ions and basal apparatus. 
\correctchoice Dynein and microtubules. 
\choice Basal apparatus and hook. 
\choice Sensory and motor neurons. 
\end{choices}

\question
Which of the following about phloem is false?

\begin{choices}
\correctchoice The plasma membrane is modified to form a sieve plate. 
\choice Sugars move by passive diffusion. 
\choice The cell has cytoplasm. 
\choice Companion cells provide proteins to sieve-tube elements. 
\choice Sieve-tube elements do not have a nucleus. 
\end{choices}
